\documentclass[preview]{standalone}

\usepackage{amsmath}
\usepackage{amssymb}
\usepackage{stellar}
\usepackage{definitions}

\begin{document}

\id{sylow-theorems-preliminaries}
\genpage

\section{Converse problems for Lagrange's theorem}

\begin{snippetexample}{lagrange-converse-counterexample-a4}{The converse of Lagrange's theorem fails}
    The alternating group \(A_4\) has order \(12\), but it has no subgroup of
    order \(6\). Indeed, such a subgroup would have index \(2\), hence would
    be normal and therefore a union of conjugacy classes. The conjugacy
    classes in \(A_4\) have cardinalities \(1,3,4,4\), with the identity in
    the singleton class. No union containing the identity has cardinality
    \(6\).
\end{snippetexample}

\begin{snippettheorem}{lagrange-converse-finite-abelian-groups}{Converse of Lagrange for finite abelian groups}
    Let \(G\) be a finite \abeliangroup of order \(n\). For every positive
    divisor \(d\) of \(n\), the group \(G\) has a \subgroup of order \(d\).
\end{snippettheorem}

\begin{snippetproof}{lagrange-converse-finite-abelian-groups-proof}{lagrange-converse-finite-abelian-groups}{Converse of Lagrange for finite abelian groups}
    First let a prime \(p\) divide \(|G|\), and prove by induction on \(|G|\)
    that \(G\) contains an element of order \(p\). The case \(|G|=p\) is
    immediate. Choose \(1_G\neq h\in G\), of order \(m\). If \(p\divides m\),
    then \(h^{m/p}\) has order \(p\). Otherwise, with \(H=\gengrp{h}\), the
    abelian quotient \(G/H\) has smaller order divisible by \(p\). By
    induction it contains an element \(gH\) of order \(p\). Then \(g^p\in H\),
    so \(g^{pm}=1_G\), while \((gH)^m\neq H\) because \(p\nmid m\). Hence
    \(g^m\neq1_G\) and \({(g^m)}^p=1_G\), so \(g^m\) has order \(p\).

    Now induct on \(d\). The case \(d=1\) is trivial. Choose a prime
    \(p\divides d\) and a subgroup \(N\) of order \(p\). Since \(G\) is
    abelian, \(N\normalsubgrp G\). The group \(G/N\) has order \(n/p\), and
    \(d/p\divides n/p\). By induction it has a subgroup \(K/N\) of order
    \(d/p\). The subgroup correspondence theorem gives
    \(|K|=|N||K/N|=d\).
\end{snippetproof}

\begin{snippettheorem}{finite-p-group-normal-subgroup-ladder}{Normal subgroup ladder in finite \(p\)-groups}
    Let \(G\) be a finite \group of order \(p^n\), where \(p\) is prime.
    For every integer \(m\) with \(0\leq m\leq n\), there exists
    \(H\normalsubgrp G\) of order \(p^m\).
\end{snippettheorem}

\begin{snippetproof}{finite-p-group-normal-subgroup-ladder-proof}{finite-p-group-normal-subgroup-ladder}{Normal subgroup ladder in finite \(p\)-groups}
    Induct on \(m\). The trivial subgroup handles \(m=0\). For \(m>0\), the
    center of the non-trivial \(p\)-group \(G\) is non-trivial. Since the
    center is abelian and its order is divisible by \(p\), the preceding
    theorem gives a subgroup \(N\subgroupleq\groupcenter(G)\) of order \(p\).
    Hence \(N\normalsubgrp G\). The quotient \(G/N\) has order \(p^{n-1}\),
    so by induction it contains a normal subgroup \(H/N\) of order
    \(p^{m-1}\). The subgroup correspondence theorem gives
    \(H\normalsubgrp G\) and \(|H|=p^m\).
\end{snippetproof}

\section{Translation action on subsets}

\begin{snippetlemma}{fixed-cardinality-subsets-action-lemma}{Action on subsets of fixed cardinality}
    Let \(G\) be a finite \group of order \(n\), let \(1\leq d\leq n\), and
    let
    \[
        X_d=\{A\subseteq G\suchthat |A|=d\}.
    \]
    The rule \(Ag=\{ag\suchthat a\in A\}\) defines a right action of \(G\)
    on \(X_d\). Every orbit has at least \(n/d\) elements. It has exactly
    \(n/d\) elements \ifandonlyif it contains a subgroup of \(G\); in that
    case it contains exactly one subgroup.
\end{snippetlemma}

\begin{snippetproof}{fixed-cardinality-subsets-action-lemma-proof}{fixed-cardinality-subsets-action-lemma}{Action on subsets of fixed cardinality}
    Right translation is bijective, so it preserves cardinality, and
    \(A1_G=A\), \((Ag)h=A(gh)\). Thus it defines the stated action.

    Let \(A\in X_d\) and choose \(a\in A\). The translate \(B=Aa^{-1}\)
    lies in the same orbit and contains \(1_G\). If
    \(S=\stabilizer_G(B)\), then every \(s\in S\) satisfies
    \(s=1_Gs\in Bs=B\), so \(S\subseteq B\). Hence \(|S|\leq d\), and
    orbit--stabilizer gives \(|BG|=n/|S|\geq n/d\). Equality holds exactly
    when \(S=B\), in which case \(B\) is a subgroup.

    Conversely, if an orbit contains a subgroup \(H\) of order \(d\), that
    orbit consists of the right cosets \(Hg\) and has \(n/d\) elements. A
    right coset \(Hg\) can be a subgroup only if it contains \(1_G\), which
    forces \(g\in H\) and \(Hg=H\). Thus the orbit contains exactly one subgroup.
\end{snippetproof}

\end{document}
